% \iffalse
%<*core>
% \fi
% \section{Tinjauan Pustaka}
% This section documents the core implementation of the \texttt{skripsi} class.
% \begin{macro}{\RequirePackage}
% Load the \textsf{kvoptions} package for key-value option handling.
%    \begin{macrocode}
\RequirePackage{kvoptions}
%    \end{macrocode}
% \end{macro}
%
% \begin{macro}{\SetupKeyvalOptions}
% Setup the key-value options system with family name and prefix.
%    \begin{macrocode}
\SetupKeyvalOptions{
    family=skripsikvopt,
    prefix=skripsikvopt@
}
%    \end{macrocode}
% \end{macro}
%
% \subsection{Boolean Options}
%
% \begin{macro}{\DeclareBoolOption}
% Declare boolean options for draft marking and title display.
% Default values: |draftmark=false|, |primarytitle=true|, |secondarytitle=true|.
%    \begin{macrocode}
\DeclareBoolOption[false]{draftmark}
\DeclareBoolOption[true]{primarytitle}
\DeclareBoolOption[true]{secondarytitle}
%    \end{macrocode}
% \end{macro}
%
% \subsection{String Options}
%
% \begin{macro}{\DeclareStringOption}
% Declare string options for section numbering depth, table of contents depth,
% and cover page style. Default values: |sectiondepth=3|, |tocdepth=3|, |cover=default|.
%    \begin{macrocode}
\DeclareStringOption[3]{sectiondepth}
\DeclareStringOption[3]{tocdepth}
\DeclareStringOption[default]{cover}
%    \end{macrocode}
% \end{macro}
%
% \subsection{Complementary Options}
%
% Define complementary options that set boolean values to false.
% |noprimarytitle| sets |primarytitle=false|, and 
% |nosecondarytitle| sets |secondarytitle=false|.
%    \begin{macrocode}
\DeclareComplementaryOption{noprimarytitle}{primarytitle}
\DeclareComplementaryOption{nosecondarytitle}{secondarytitle}
%    \end{macrocode}
%
% Process all declared key-value options.
%    \begin{macrocode}
\ProcessKeyvalOptions*
%    \end{macrocode}
%
% \begin{macro}{\LoadClass}
% Load the base \textsf{book} class with A4 paper, 12pt font, and one-sided layout.
%    \begin{macrocode}
\LoadClass[a4paper,12pt,oneside]{book}
%    \end{macrocode}
% \end{macro}
%
% \section{Package Dependencies}
%Packages loaded by the class for layouts, fonts, lists, and references.
%    \begin{macrocode}
\RequirePackage{environ}
\RequirePackage{polyglossia}
\RequirePackage{fontspec}
\RequirePackage{titling}
\RequirePackage{fancyhdr}
\RequirePackage{geometry}
\RequirePackage{setspace}
\RequirePackage{graphicx}
\RequirePackage{tikz}
\RequirePackage{tikzpagenodes}
\usetikzlibrary{calc}
\RequirePackage{svg}
\RequirePackage{etoolbox}
\RequirePackage{tocbibind}
\RequirePackage{indentfirst}
\RequirePackage{float}
\RequirePackage{atbegshi}
\RequirePackage{amsmath,amsfonts,amssymb}
\RequirePackage{hyperref}
\RequirePackage{hyperxmp}
\RequirePackage[acronym, nogroupskip, toc]{glossaries}
%    \end{macrocode}
%    
%\begin{macro}{Language defaults}
%Set main and other languages.
%    \begin{macrocode}
\setmainlanguage[variant=indonesian]{malay}
\setotherlanguage{english}
%    \end{macrocode}
%\end{macro}
%
%\begin{macro}{Draft watermark}
%Conditionally add a \textsf{draftwatermark} and line numbers.
%    \begin{macrocode}
\ifskripsikvopt@draftmark
  \RequirePackage{draftwatermark}
  \DraftwatermarkOptions{%
    text={\sffamily\bfseries
      DRAFT 
      \href{https://github.com/nuluh/thesis}{%
        \color{magenta}\fbox{\bfseries\normalsize\texttt{ver.\ 2025.j5}}%
      }%
      \ \normalsize{Copyright~\textcopyright~Rifqi D. Panuluh}%
    },
    color={[gray]{0.8}},
    fontsize=1.5em,
    angle=90,
    hpos=20cm,
    vpos=20cm
  }
  \RequirePackage[left]{lineno}
  \linenumbers
\fi
%    \end{macrocode}
%\end{macro}
%
% 
% \begin{macro}{\skripsi@setup@layout}
% Applies default margins and spacing for the thesis class.
%    \begin{macrocode}
\newcommand*\skripsi@setup@layout{%
  \geometry{left=4cm, top=3cm, right=3cm, bottom=3cm}%
  \setlength{\parskip}{0.5em}%
}
%    \end{macrocode}
% \end{macro}
% call it during class load
%    \begin{macrocode}
\skripsi@setup@layout
%    \end{macrocode}
%
%\begin{macro}{Fonts}
%Set fonts with \textsf{fontspec}.
%    \begin{macrocode}
\defaultfontfeatures{Ligatures=TeX}
\setmainfont{Times New Roman}[
  SmallCapsFont  = {Latin Modern Roman}, % fallback for \textsc
]
\setsansfont{Arial}
\setmonofont{Courier New}
%    \end{macrocode}
%\end{macro}
%
%\begin{macro}{Page styles}
%Fancyhdr setups for main text and front matter.
%    \begin{macrocode}
\fancypagestyle{fancy}{%
    \fancyhf{}
    \fancyhead[R]{\nouppercase{\rightmark}}
    \fancyhead[L]{\nouppercase{\leftmark}}
    \fancyfoot[C]{\thepage}
}
\fancypagestyle{fancyplainfrontmatter}{%
    \renewcommand{\headrulewidth}{0pt}
    \fancyfoot[C]{\thepage}
}
\fancypagestyle{fancyplain}{%
    \fancyhf{}
    \setlength{\headheight}{15pt} % avoid fancyhdr warning
    \renewcommand{\headrulewidth}{0pt}
    \fancyhead[R]{\thepage}
}
%    \end{macrocode}
%\end{macro}
% 
% define frontmatter hooks for page numbering name
%    \begin{macrocode}
\makeatletter
%    \end{macrocode}
%\begin{macro}{\mainmatter,\backmatter}
%Back matter uses roman numerals while continuing the page counter from main matter.
%    \begin{macrocode}
\newcounter{savepagenumber}
\renewcommand\frontmatter{%
  \if@openright
    \cleardoublepage
  \else
    \clearpage
  \fi
  \skripsi@frontmatter@pagenumbering
  \skripsi@frontmatter@pagestyle
  \@mainmatterfalse
}
\renewcommand\mainmatter{%
  \cleardoublepage
  \setcounter{savepagenumber}{\value{page}}
  \@mainmattertrue
  \skripsi@mainmatter@pagenumbering
  \skripsi@mainmatter@pagestyle
}
\renewcommand\backmatter{%
  \if@openright
    \cleardoublepage
  \else
    \clearpage
  \fi
  \skripsi@backmatter@pagenumbering
  \skripsi@backmatter@pagestyle
  \setcounter{page}{\value{savepagenumber}}%
  \@mainmatterfalse
}
%    \end{macrocode}
%\end{macro}
%
% \section{Sectioning Style}
% Apply heading layout and numbering depth.
%    \begin{macrocode}
\RequirePackage{titlesec}
%    \end{macrocode}
% \subsection{Chapter}
%Chapter title is formatted as shown below to "BAB I", "BAB II", etc. where \verb|\chaptername| is defined in the polyglossia language settings,
%    \begin{macrocode}
\titlespacing{\chapter}{0pt}{0cm}{*1.5}
\titleformat{\chapter}[display]
  {\normalsize\bfseries\centering}
  {\MakeUppercase\chaptername~\Roman{chapter}}
  {1ex}
%    \end{macrocode}
% followed by chapter title text itself which uppercased.
%    \begin{macrocode}
  {\MakeUppercase}
%    \end{macrocode}
% \subsection{Section and Subsection}
%The following line defines section and subsection formatting to normal size, bold font, followed with their respective numbering then horizontal separation between the numbering label and the title text.
%    \begin{macrocode}
\titleformat{\section}
    {\normalsize\bfseries}{\thesection}{1em}{}
\titleformat{\subsection}
    {\normalsize\bfseries}{\thesubsection}{1em}{}
%    \end{macrocode}
%    Set section numbering depth to include up to \texttt{subsubsection} (1.1.1.1 Subsubsection).
%    \begin{macrocode}
\setcounter{secnumdepth}{3}
%    \end{macrocode}
% \subsection{Paragraph}
% Define indent lengths for section, subsection, and subsubsection levels.
%    \begin{macrocode}
\newlength{\skripsi@section@indent}
\newlength{\skripsi@subsection@indent}
\newlength{\skripsi@subsubsection@indent}
%    \end{macrocode}
% Set default indent lengths to zero.
%    \begin{macrocode}
\setlength{\skripsi@section@indent}{\z@}
\setlength{\skripsi@subsection@indent}{\z@}
\setlength{\skripsi@subsubsection@indent}{\z@}
%    \end{macrocode}
% Then, prepend each sectioning command to set the left skip according to the defined indent lengths with ensure paragraph break before setting the indent.
%    \begin{macrocode}
\pretocmd{\section}{%
  \par
  \setlength{\leftskip}{\skripsi@section@indent}%
}{}{}

\pretocmd{\subsection}{%
  \par
  \setlength{\leftskip}{\skripsi@subsection@indent}%
}{}{}

\pretocmd{\subsubsection}{%
  \par
  \setlength{\leftskip}{\skripsi@subsubsection@indent}%
}{}{}
%    \end{macrocode}
%\section{Cover Page}
%\begin{macro}{\maketitle}
% Redefine \cs{maketitle} to support multiple cover page styles with fallback options.
% The macro first saves the original \cs{maketitle} from the book class, then provides
% a flexible system that checks for user-defined cover pages in the \texttt{frontmatter/}
% directory, with embedded fallback templates if custom files are not found.
%
% \subsection{Primary Title Page}
%
% Save the original \cs{maketitle} command before redefining it.
%    \begin{macrocode}
\let\original@maketitle\maketitle

\renewcommand{\maketitle}{%
%    \end{macrocode}
%
% If the |primarytitle| option is enabled, generate the primary cover page.
% Three modes are supported based on the |cover| option:
% \begin{itemize}
%   \item |cover=legacy|: Uses the original book class \cs{maketitle}
%   \item Custom file: Uses \texttt{frontmatter/maketitle.tex} if it exists
%   \item Embedded: Uses the built-in template as fallback
% \end{itemize}
%    \begin{macrocode}
  \ifskripsikvopt@primarytitle
    \ifdefstring{\skripsikvopt@cover}{legacy}{%
      \original@maketitle
      \IfFileExists{frontmatter/maketitle.tex}{%
      \ClassWarning{skripsi.cls}{Omit 'cover=legacy' option to use 
                    user-defined primary cover page 
                    'frontmatter/maketitle.tex'}
      }{}%
    }{%
      \IfFileExists{frontmatter/maketitle.tex}{%
        \begin{titlepage}

\centering
{\large\selectfont \textbf{\MakeUppercase{\thetitle}}}
\vspace{1.5cm}

{\textbf{Tugas Akhir}\par}
{\textbf{%
diajukan untuk memenuhi salah satu syarat memperoleh gelar sarjana \\
dari \theprogram, \\
\thefaculty \\
\theuniversity}\par}

\vfill

{\textbf{\theauthor}\par}
{\textbf{\theauthorid}\par}

\vspace{\baselineskip}

\includegraphics[width=3.6cm]{frontmatter/img/logo.png}

\vfill

{\large\selectfont
\textbf{\theprogram} \par
\textbf{\thefaculty} \par
\textbf{\theuniversity} \par
\textbf{\thecity} \par
\textbf{\theyearofsubmission} \par}

\end{titlepage}%%
        \ClassWarning{skripsi.cls}{Using user-defined cover page 
                      'frontmatter/maketitle.tex'}%
      }{%
%    \end{macrocode}
%
% \subsubsection{Embedded Primary Cover Template}
% If no custom file exists, use the embedded template with standard formatting:
% centered layout, university logo, author information, and program details. But checking required metadata: |\title|, |\author|, |\date|, etc. first before proceeding. Due to the use of titling package, error handling is done different for title, author, and date.
%    \begin{macrocode}
      \ifdefined\thetitle\else \ClassError{skripsi.cls}{Missing \protect\title}{}\fi
      \ifdefined\theauthor\else \ClassError{skripsi.cls}{Missing \protect\author}{}\fi
      \ifdefined\thedate\else \ClassError{skripsi.cls}{Missing \protect\date}{}\fi
      \@check{program}
      \@check{faculty}
      \@check{university}
      \@check{yearofsubmission}
%    \end{macrocode}
% And then proceed to generate the cover page.
%    \begin{macrocode}
        {%
          \begin{titlepage}
          \newgeometry{margin=3cm}%
          \noindent
          \begin{minipage}[t][0.49\textheight][t]{\textwidth}
          \centering
          {\fontsize{14pt}{16pt}\selectfont \textbf{\MakeUppercase{Tugas Akhir}}\par}
          \vspace{1.5cm}
          {\fontsize{14pt}{16pt}\selectfont \textbf{\MakeUppercase{\thetitle}}\par}
          \vfill
          \end{minipage}
          \begin{tikzpicture}[remember picture, overlay]
              % This respects your geometry margins
              \coordinate (center) at (current page text area.center);
              
              % The logo at center of page (ANCHOR POINT) shift up by 1cm
              \node at ($(center) + (0cm,1cm)$)
              {\includegraphics[width=5cm]{frontmatter/img/logo.png}};
          \end{tikzpicture}%
          \begin{minipage}[b][0.49\textheight][b]{\textwidth}
          \centering

          \vfill

          \textbf{Disusun oleh:} \\
          {\fontsize{14pt}{16pt}\selectfont \textbf{\theauthor}} \\
          {\fontsize{14pt}{16pt}\selectfont \textbf{\theauthorid}} \\

          \vfill

          {\fontsize{12pt}{14pt}\selectfont
          \textbf{\MakeUppercase\theprogram} \\
          \textbf{\MakeUppercase\thefaculty} \\
          \textbf{\MakeUppercase\theuniversity} \\
          \textbf{\theyearofsubmission}
          }
          \end{minipage}
          \end{titlepage}%
        }
      }%
    }%
  \else
%    \end{macrocode}
%
% If |primarytitle| is disabled but a custom file exists, warn the user.
%    \begin{macrocode}
      \IfFileExists{frontmatter/maketitle.tex}{%
        \ClassWarning{skripsi.cls}{Omit 'noprimarytitle' option to use 
                      user-defined cover page 
                      'frontmatter/maketitle.tex'}%
      }{}%
  \fi
%    \end{macrocode}
%
% \subsection{Secondary Title Page}
%
% If the |secondarytitle| option is enabled, generate the secondary cover page
% after a page break. Similar to the primary page, it checks for a custom file
% in \texttt{frontmatter/maketitle\_secondary.tex} and falls back to an
% embedded template if not found.
%    \begin{macrocode}
  \ifskripsikvopt@secondarytitle
    \clearpage
    \IfFileExists{frontmatter/maketitle_secondary.tex}{%
      \input{frontmatter/maketitle_secondary}%
      \ClassWarning{skripsi.cls}{Using user-defined secondary cover page}
    }{%
%    \end{macrocode}
%
% \subsubsection{Embedded Secondary Cover Template}
%
% The embedded secondary template includes additional submission purpose text
% and maintains similar styling to the primary cover. %TODO: secondary cover paragraph text
%    \begin{macrocode}
      {%
        \begin{titlepage}
        \newgeometry{margin=3cm, left=4cm}
        \noindent
        \begin{minipage}[t][0.49\textheight][t]{\textwidth}
        \centering
        {\fontsize{14pt}{16pt}\selectfont \textbf{\MakeUppercase{Tugas Akhir}}\par}
        \vspace{1.5cm}
        {\fontsize{14pt}{16pt}\selectfont \textbf{\MakeUppercase{\thetitle}}\par}
        \vspace{1cm}
        {\normalsize\selectfont \thesubmissionstatement\par}
        \vfill
        \end{minipage}
        \begin{tikzpicture}[remember picture, overlay]
            % This respects your geometry margins
            \coordinate (center) at (current page text area.center);
            
            % The logo at center of page (ANCHOR POINT) shift up by 1cm
            \node at ($(center) + (0cm,1cm)$)
                {\includegraphics[width=5cm]{frontmatter/img/logo.png}};
        \end{tikzpicture}%
        \begin{minipage}[b][0.49\textheight][b]{\textwidth}
        \centering
        \vfill

        \textbf{Disusun oleh:} \\
        {\fontsize{14pt}{16pt}\selectfont \textbf{\theauthor}} \\
        {\fontsize{14pt}{16pt}\selectfont \textbf{\theauthorid}} \\
        \vfill
        {\fontsize{12pt}{14pt}\selectfont
        \textbf{\MakeUppercase\theprogram} \\
        \textbf{\MakeUppercase\thefaculty} \\
        \textbf{\MakeUppercase\theuniversity} \\
        \textbf{\theyearofsubmission}
        }
        \end{minipage}
        \end{titlepage}%
      }
      \ClassWarning{skripsi.cls}{File frontmatter/maketitle_secondary.tex 
                    not found, using embedded secondary cover page}%
    }%
  \else
%    \end{macrocode}
% 
% If |secondarytitle| is disabled but a custom file exists, warn the user.
%    \begin{macrocode}
      \IfFileExists{frontmatter/maketitle_secondary.tex}{%
        \ClassWarning{skripsi.cls}{Omit 'nosecondarytitle' option to use 
                      user-defined secondary cover page 
                      'frontmatter/maketitle_secondary.tex'}%
      }{}%
  \fi
}
%    \end{macrocode}
%\end{macro}
%
%\section{Table of Contents, List of Figures, and List of Tables}
%    \begin{macrocode}
\RequirePackage[titles]{tocloft}
%    \end{macrocode}
% \subsection{Chapter}
% Redefine printed chapter font in the table of contents to be UPPERCASED while keeping the actual string and pdf bookmark table of contents in Title Case.
%  First, turns chapter name "Bab" to be UPPERCASED ("BAB").
%    \begin{macrocode}
\renewcommand{\cftchappresnum}{\MakeUppercase\chaptername~}
%    \end{macrocode}
%    Then followed with horizontal space then the chapter number.
%    \begin{macrocode}
\renewcommand{\cftchapaftersnum}{\quad}
%    \end{macrocode}
%    Finally, followed by the chapter title which also UPPERCASED.
%    \begin{macrocode}
\renewcommand{\cftchapfont}{\normalsize\MakeUppercase}
%    \end{macrocode}
% Following line is necessary to fix the error (thanks to \href{https://latex.org/forum/viewtopic.php?p=83118&sid=a22ccb456531faabf65070c18af69680#p83118}{Latex.org} for elegant solution). This is elegant approach since it preserves the actual string casing that is passed into |\chapter|, so it appears correctly in PDF bookmarks while displaying UPPERCASE in the printed table of contents.
%    \begin{macrocode}
\patchcmd{\l@chapter}
  {\cftchapfont #1}% original
  {\cftchapfont{#1}}% patched
  {}{}            % success/failure handlers
%    \end{macrocode}
% \subsection{TOC, LOF, and LOT Length Settings}
% See "Figure 1 Layout of a ToC (LoF, LoT)" in tocloft documentation entry for details.
%    \begin{macrocode}
\setlength{\cftbeforetoctitleskip}{0cm}
\setlength{\cftbeforeloftitleskip}{0cm}
\setlength{\cftbeforelottitleskip}{0cm}
\setlength{\cftbeforechapskip}{0em}
\setlength{\cftchapindent}{0pt}
\setlength{\cftsecindent}{0em}
\setlength{\cftsubsecindent}{2em}
\setlength{\cftchapnumwidth}{3.5em}
\setlength{\cftsecnumwidth}{2em}
\setlength{\cftsubsecnumwidth}{2.5em}
\setlength{\cftfignumwidth}{6em}
\setlength{\cfttabnumwidth}{4em}
\renewcommand \cftchapdotsep{1}
\renewcommand \cftsecdotsep{1}
\renewcommand \cftsubsecdotsep{1} 
\renewcommand \cftfigdotsep{1.5}
\renewcommand \cfttabdotsep{1.5}
\renewcommand{\cftchapleader}{\normalfont\cftdotfill{\cftsecdotsep}}
\renewcommand{\cftchappagefont}{\normalfont}
%    \end{macrocode}
% \subsection{Figure and Table Numbering}
% Redefine chapter and section numbering formats.
%    \begin{macrocode}
\renewcommand{\thechapter}{\Roman{chapter}}
\renewcommand\thesection{\arabic{chapter}.\arabic{section}}
%    \end{macrocode}
% Redefine figure, table, and equation numbering to include chapter number as prefix. (e.g., Figure 1.1, Table 2.3, Equation 3.5)
%    \begin{macrocode}
\renewcommand{\thefigure}{\arabic{chapter}.\arabic{figure}}
\renewcommand{\thetable}{\arabic{chapter}.\arabic{table}}
\renewcommand{\theequation}{\arabic{chapter}.\arabic{equation}}
%    \end{macrocode}
% Add "Gambar" and "Tabel" prefixes before figure and table numbers in the list of figures and list of tables.
%    \begin{macrocode}
\renewcommand{\cftfigpresnum}{\figurename~}
\renewcommand{\cfttabpresnum}{\tablename~}
%    \end{macrocode}
% Redefine TOC, LOF, and LOT page titles to be centered, bold, and UPPERCASED. By giving \cs{hfill} before and after the title font command, the title is centered. This doesnt affects PDF bookmark entries, which remain in Title Case.
%    \begin{macrocode}
\renewcommand{\cfttoctitlefont}{\hfill\bfseries\MakeUppercase}
\renewcommand{\cftaftertoctitle}{\hfill} % https://tex.stackexchange.com/a/255699/394075
\renewcommand{\cftloftitlefont}{\hfill\bfseries\MakeUppercase}
\renewcommand{\cftafterloftitle}{\hfill}
\renewcommand{\cftlottitlefont}{\hfill\bfseries\MakeUppercase}
%    \end{macrocode}
% Set section numbering depth in the table of contents page according to the |tocdepth| option. This is different from PDF bookmark depth which is controlled by \textsf{hyperref} package. To set PDF bookmark depth, see |bookmarks| option in \textsf{hyperref}.
%    \begin{macrocode}
\setcounter{tocdepth}{\skripsikvopt@tocdepth}
%    \end{macrocode}
% \section{Abstract}
% \subsection{Keywords}
% \label{sec:keyword-macro}
% \subsubsection{Engine-Agnostic String Comparison}
%
% Define a wrapper macro for string comparison that works across different \LaTeX{}
% engines without requiring external packages. The |\pdf@strcmp| macro provides a
% unified interface to engine-specific string comparison primitives.
%
%    \begin{macrocode}
\makeatletter
\ifluatex
  \let\pdf@strcmp\pdfstrcmp
\else\ifxetex
  \let\pdf@strcmp\strcmp
\fi\fi
%    \end{macrocode}
% \begin{macro}{\keywords}
% \changes{v1.0}{2025/11/09}{Initial version}
% The keyword system provides automatic alphabetical sorting of comma-separated
% keywords and integrates them into both the document display and PDF metadata.
% This implementation requires Lua\LaTeX{} compilation.
% 
% Sort a comma-separated list of keywords alphabetically.
% This serves as the \LaTeX{} interface to the Lua |sort_keywords| function
% defined in Section~\ref{sec:lua-keyword-sort}.
% 
% \paragraph{Implementation Strategy}
% This implementation delegates both parsing and sorting entirely to Lua:
% \begin{enumerate}
% \item Pass the raw comma-separated string to Lua via |\directlua|
% \item Lua handles delimiter splitting, whitespace trimming, and sorting
% \item Return the sorted, formatted string to \LaTeX{}
% \end{enumerate}
% 
% This approach is simpler and more maintainable than traditional \TeX{} string
% manipulation, avoiding catcode issues and complex loop constructs. The resulting
% sorted list is stored in |\sortedlist| for use in the |\keywords| macro.
% 
% Pass the entire keyword string to |sort_keywords| and store the result in
% |\sortedlist|. The |\luaescapestring| command ensures that special characters
% (backslashes, quotes, etc.) are properly escaped for Lua parsing, while
% |\unexpanded| prevents premature macro expansion.
% Display sorted keywords and embed them in PDF metadata.
% This macro provides the user-facing interface for keyword input, handling both
% visual presentation and metadata generation.
%
% ^^A TODO: Add capability to preserve basic formatting when sorting like |\textit|, |\textbf| in keywords.
%    \begin{macrocode}
\newcommand{\keywords}[1]{%
\par\bigskip\noindent
    \ifLuaTeX
        \edef\sortedlist{%
            \directlua{
            local input = "\luaescapestring{\unexpanded{#1}}"
            local sorted = sort_keywords(input)
            tex.sprint(sorted)
            }%
        }%
%    \end{macrocode}
% If not compiled with Lua\LaTeX, issue a class warning to consider using Lua\LaTeX for this sorting.
%    \begin{macrocode}
    \else
        \ClassWarning{skripsi}{%
            Keyword sorting requires LuaLaTeX; 
            consider compiling with LuaLaTeX for auto-sort keywords%
        }%
        \gdef\sortedlist{#1}%
    \fi
%    \end{macrocode}
% \paragraph{Multilingual Support}
% The macro automatically detects the document language using |\languagename| and
% adjusts the keyword label accordingly:
% \begin{itemize}
% \item Malay (|malay|): displays as ``Kata Kunci:'' in upright text
% \item Other languages: displays as ``\textit{Keywords:}'' in italic text
% \end{itemize}
%    \begin{macrocode}
    \ifnum\pdf@strcmp{\languagename}{malay}=0
        Kata Kunci: \sortedlist%
        \else
            \textit{Keywords: \sortedlist}%
        \fi
%    \end{macrocode}
%    
% \paragraph{PDF Metadata Integration}
% Sorted keywords are embedded in the PDF document information dictionary through
% three mechanisms to ensure maximum compatibility:
% \begin{enumerate}
% \item Direct embedding via |\special{pdf:docinfo}| for Xe\LaTeX/xdvipdfmx workflows
% \item Updating hyperref's internal metadata via |\hypersetup| (if loaded)
% \item Proper escaping of special characters using |\pdfstringdef| to ensure
%       PDF string compliance
% \end{enumerate}
%    \begin{macrocode}
  \gdef\SortedKeywords{\sortedlist} \pdfstringdef\kwpdf{\SortedKeywords}%
%    \end{macrocode}
%    For XeLaTeX and xdvipdfmx workflows, directly embed the sanitized keywords into the PDF's document information dictionary. The last invocation of this special takes precedence.
%    \begin{macrocode}
  \immediate\special{pdf:docinfo << /Keywords (\kwpdf) >>}%
%    \end{macrocode}
%    Attempt to update hyperref's internal PDF metadata as well (harmless if hyperref is unavailable or already finalized).
%    \begin{macrocode}
  \hypersetup{pdfkeywords={\SortedKeywords}}%
}
%    \end{macrocode}
% \end{macro}
% \subsection{Abstract Word Count Warning}
% \begin{environment}{abstract}
% \changes{v1.0}{2025/11/09}{Initial version}
% 
% Define the abstract environment with word count validation.
% This implementation uses \textsf{environ} to capture the entire environment body,
% enabling word count analysis before typesetting.
%
% Due to the inherent limitations of pure \LaTeX{} for text processing and parsing,
% word counting is implemented only for Lua\LaTeX{} workflows. When compiled with
% Lua\LaTeX, the environment captures the abstract body and passes it to a Lua
% function (|count_words|) for accurate word counting. If the word count exceeds
% the threshold defined in |\maxabstractlengthwarning|, a class warning is issued
% to both the log file and console output.
%
% \begin{itemize}
% \item Word counting is unavailable in pdf\LaTeX{} and \XeLaTeX{} workflows
% \item Macro expansion within the abstract may affect word count accuracy
% \item Special characters and mathematical content may be counted unexpectedly
% \end{itemize}
%
%    \begin{macrocode}
\newcommand\abstractname{Abstract}
\newcommand{\maxabstractlengthwarning}{200} % TODO: make class warning if user trying to redefine this and consider to use LuaLaTeX for word count
\NewEnviron{abstract}{%
%    \end{macrocode}
% Store the environment body globally for potential reuse (e.g., in PDF metadata).
%    \begin{macrocode}
  \global\let\theabstract\BODY
%    \end{macrocode}
% Perform word count validation only when compiled with Lua\LaTeX.
%    \begin{macrocode}
  \ifLuaTeX
%    \end{macrocode}
% Execute the Lua word counting function and store the result in |\abstractwordcount|.
% The |\luaescapestring| and |\unexpanded\expandafter| combination ensures that
% special characters and macros in the abstract are properly escaped before being
% passed to Lua, preventing compilation errors.
%    \begin{macrocode}
    \edef\abstractwordcount{%
      \directlua{
        local text = "\luaescapestring{\unexpanded\expandafter{\BODY}}"
        local count = count_words(text)
        tex.print(count)
      }%
    }%
%    \end{macrocode}
% Compare the word count against the maximum threshold. If exceeded, issue a
% class warning that appears in both the |.log| file and console output.
%    \begin{macrocode}
    \ifnum\abstractwordcount>\maxabstractlengthwarning\relax
      \ClassWarning{skripsi}{%
        Abstract exceeds \maxabstractlengthwarning\space words 
        (current count: \abstractwordcount\space words)%
      }%
      \typeout{WARNING: Abstract exceeds \maxabstractlengthwarning\space words!}%
      \typeout{Word count: \abstractwordcount}%
    \fi
  \fi
%    \end{macrocode}
% Typeset the abstract as a chapter with the appropriate heading.
%    \begin{macrocode}
  \chapter{\abstractname}
  \BODY
}
\makeatother
%    \end{macrocode}
% \end{environment}
% \iffalse
%</core>
% \fi